% displayname : Rayonnement thermique
\documentclass[12pt,fancy]{cours}
\Chapit{2}{3}
\begin{document}
\maketitre[***]{
Ce chapitre est un complément sur la notion de rayonnement et d'équilibre thermique d'un corps. Nous adoptons une approche descriptive, aucune démonstration ne sera donnée, ni aucune formule à connaître par cœur. Un exemple important de bilan radiatif concerne la planète Terre et son atmosphère, pour laquelle nous proposons plusieurs modèles afin d'expliquer la température moyenne de surface.
}
{
% \item Cours Thermodynamique 1 à 3 - 0
\item Documents de cours Thermodynamique 2 - 1
% \item TD Thermodynamique 1 à 3
}


%%%%%%%%%%%%%%%%%%%%%%%%%%%%%%%%%%%%%%%%%%%%%%%%%%%%%%%%%%%%%%%%%%%%%%%%%%
\section{Rayonnement  d'équilibre thermique}
\subsection{Définitions}

\parag{Rayonnement d'équilibre thermique}
La \imp{radiation} est un transfert thermique qui s'effectue par émission ou absorption d'un rayonnement électromagnétique. Ce transfert peut s'effectuer dans le vide, puisque les ondes électromagnétiques n'ont pas besoin de milieu matériel pour se propager.

Lorsqu'un métal est chauffé par le forgeron, il émet une lumière rougeâtre qui devient progressivement orangée puis jaune lorsque sa température augmente. De façon générale, tout corps émet un rayonnement électromagnétique suite à l'agitation thermique des atomes, appelé rayonnement thermique.  Lorsque les photons qui constituent le rayonnement thermique émis par les atomes du corps émetteur de température $T$ sont autant absorbés que réémis par ce corps, on a atteint une situation d'équilibre.
\begin{definition}
Le \imp{rayonnement d'équilibre thermique} est le rayonnement thermique émis par un corps de température $T$ à l'équilibre avec le rayonnement qu'il émet.
\end{definition}
\noindent
Le rayonnement d'équilibre thermique dépend uniquement de $T$, mais pas des détails du processus d'émission et d'absorption, et donc du corps émetteur particulier.


%Lorsqu'un corps est soumis à un rayonnement électromagnétique, l'énergie électromagnétique incidente peut être \important{absorbée}, \important{transmise} (elle traverse le corps), ou \important{réfléchie}. La fraction du rayonnement réfléchi sur le rayonnement incident est appelée l'\important{albédo}. Il s'agit donc d'un coefficient compris entre $0$ et $1$. Par exemple un albédo de $\num{0,3}$ signifie que le corps en question réfléchit $\SI{30}{\percent}$ de l'énergie incidente.
%
%\begin{figure}[h!]
%\centering
%\includegraphics[width=0.8\textwidth]{schema_energie_incidente.pdf}
%\caption{Lorsqu'un corps quelconque (dit gris) est à l'équilibre thermique, l'énergie incidente se divise en une partie réfléchie, une absorbée et une transmise. Pour un corps noir, toute l'énergie incidente est absorbée.}
%\end{figure}
%
%Une idéalisation très utile est celle du \important{corps noir}. Il s'agit d'un corps qui absorbe la totalité du rayonnement incident. Cette énergie est convertie sous forme de chaleur. À l'équilibre le corps noir atteint une certaine température qui seule permet de calculer le flux d'énergie qui sera réémis (loi de Stefan, voir ci-dessous).


% Supposons maintenant que le rayonnement incident est un rayonnement d'équilibre thermique. Alors $\varphi_i = \varphi^0(T)$.
% Si le corps est de plus à l'équilibre radiatif avec ce rayonnement d'équilibre thermique, on a aussi que $\varphi_i = \varphi_s$.
% \begin{prop}\label{prop:}
% Un corps soumis à un rayonnement incident est à l'\imp{équilibre thermodynamique} s'il est en équilibre radiatif et que le rayonnement incident est à l'équilibre thermique. Alors
% \begin{equation*}
% \varphi_i = \varphi_s = \varphi^0(T).
% \end{equation*}
% \end{prop}
% Par contre, le flux émis $\varphi_e=\varphi_a$ dépend de la nature du corps en général, et n'est donc pas un rayonnement d'équilibre thermique. Pour modéliser simplement le rayonnement émis, on utilise le modèle du corps noir.

La plupart des corps n'émettent pas un rayonnement d'équilibre thermique. On définit alors un corps modèle, dont le rayonnement émis est un rayonnement d'équilibre thermique, appelé \imp{corps noir}. Pour vérifier cette propriété, un corps noir doit absorber intégralement le rayonnement incident.
\begin{definition}\label{def:corps_noir}
Un \imp{corps noir} est un corps qui absorbe intégralement le rayonnement incident, et émet un rayonnement d'équilibre thermique.
\end{definition}
% Un corps noir vérifie donc $\varphi_i = \varphi_a$ et $\varphi_s = \varphi_e$.
\begin{remarque}
Un corps noir n'est pas de couleur noire ! Il absorbe tout le rayonnement incident. Par conséquent, la couleur perçue ne dépend pas du rayonnement incident mais seulement du rayonnement émis, et donc de la température du corps noir.
\end{remarque}

%À température ambiante (\SI{300}{K}), la loi de Wien indique qu'effectivement la couleur observée est noire. Mais au delà de \SI{2000}{K}, la loi de Wien donne une longueur d'onde maximale qui est dans le visible, par exemple un corps qui émet à une température de surface de \SI{6000}{K} comme le Soleil apparaît jaune. Et pourtant, cela n'empêche pas de modéliser le Soleil par un corps noir.
%
% Pour réaliser un corps noir en pratique, on peut prendre une enceinte opaque chauffé percée d'un petit trou par lequel s'échappe une quantité négligeable de rayonnement. Imaginons un corps noir a l'équilibre radiatif avec un rayonnement d'équilibre thermique. Alors $\varphi_i = \varphi_s = \varphi^0(T)$ d'après l'équilibre thermodynamique, et  $\varphi_i = \varphi_a$ et $\varphi_s = \varphi_e$ d'après les propriétés du corps noir.
% \begin{prop}\label{prop:}
% Un corps noir à l'équilibre thermodynamique émet un rayonnement d'équilibre thermique, soit
% $$\varphi_e=\varphi_a=\varphi^0(T).$$
% \end{prop}



\subsection{Propriétés du rayonnement d'équilibre thermique}


\parag{Loi de Planck} Le rayonnement d'équilibre thermique est caractérisé par la \imp{puissance surfacique} d'équilibre thermique $\varphi_0(T)$ (en $\SI{}{W.m^{-2}}$), fonction uniquement de la température $T$ du rayonnement. Il s'agit de la puissance par unité de surface projetée transportée par le champ électromagnétique à un corps.
\begin{remarque}
  Nous verrons plus tard que cette grandeur correspond à la projection du vecteur de Poynting $\Poy$ avec la normale  $\Vec{n}$ de la surface.
\end{remarque}
On peut décomposer la puissance surfacique en un ensemble de composantes de longueur d'onde $\lambda$, grandeur que l'on appelle puissance surfacique spectrale $\dd{\vphi_0}{\lambda}$. Max Planck a déterminé en 1900 l'expression de la puissance surfacique spectrale, démonstration au cours de laquelle il a introduit la fameuse constante~$h$ qui porte maintenant son nom. La loi de Planck s'écrit
\begin{equation}\label{eq:planck}
\begin{split}
\dd{\vphi_0}{\lambda}(\lambda,T)=\dfrac{2\pi hc^2}{\lambda^{5}}\dfrac{1}{\exp\!\left(\frac{h c}{\lambda k_{\rm{B}}T}\right)-1},
\end{split}
\end{equation}
avec $h=\SI{6,63e-34}{J.s}$ la constante de Planck et $k_{\rm{B}}=\SI{1,38e-23}{J.K^{-1}}$ celle de Boltzmann. L'allure du graphe de la puissance surfacique spectrale en fonction de la longueur d'onde $\lambda$ pour différentes températures $T$ est montrée en Fig.~\ref{fig:corps_noir}. L'\Eq{planck} n'est pas à connaître, mais elle conduit à deux lois qui elles sont à maîtriser.
% C'est l'équivalent de la densité de courant thermique $\Vec{\jth} \cdot \Vec{n}$ de la conduction thermique.
% En assimilant le rayonnement à un gaz de photons homogène, on peut montrer que ces deux grandeurs sont reliées par
% \begin{equation}
% \varphi_0(T) = \dfrac{c}{4}u_0(T),
% \end{equation}
% où $c$ est la célérité de la lumière dans le vide.
%
% Le rayonnement d'équilibre thermique comprend tout un \imp{continuum de longueurs d'onde}. Il est donc pertinent de connaître la répartition en longueurs d'onde de la densité d'énergie électromagnétique ou du flux surfacique. On définit ainsi :
% \begin{liste}
% \item la densité spectrale d'énergie volumique  $u^0_\lambda(\lambda, T)$ (en $\SI{}{J.m^{-4}}$), telle que la quantité d'énergie volumique dans l'intervalle $[\lambda,\lambda+\d \lambda]$ soit $\d u^0 = u^0_\lambda(\lambda, T)\d \lambda$ ;
% \item Le flux surfacique spectral $\varphi^0_\lambda(\lambda,T)$ (en $\SI{}{W.m^{-3}}$), tel que le flux surfacique dans l'intervalle $[\lambda,\lambda+\d \lambda]$ soit $\d \varphi^0 = \varphi^0_\lambda(\lambda, T)\d \lambda$.
% \end{liste}

\begin{figure}[H]
	\centering
	\includegraphics[scale=1.5]{black_Body.pdf}
	\caption{Spectre d'émission prédit par la loi de Planck. La puissance surfacique spectrale d'équilibre thermique est représentée en fonction de la longueur d'onde $\lambda$ pour trois températures $T=\SI{3000}{K}$, $\SI{4000}{K}$ et $\SI{5000}{K}$. La bande colorée entre $\SI{400}{nm}$ et $\SI{750}{nm}$ reprérsente la partie visible du spectre. La loi de déplacement de Wien donne la position du pic d'émission $\lambda_{\rm{max}}(T)$ (courbe verte tiretée). L'aire sous chaque courbe croît avec la température comme $T^4$ suivant la loi de Stefan.}
	\label{fig:corps_noir}
\end{figure}


\parag{Loi de déplacement de Wien}
La fonction $\lambda \mapsto \dd{\vphi_0}{\lambda}(\lambda,T)$ passe par un maximum unique pour la longueur d'onde $\lambda_{\rm max}(T)$ telle que :
\begin{theorem}[Loi de déplacement de Wien]\label{form:wien}
\begin{equation*}
\lambda_{\rm{max}}\cdot T \simeq \SI{2900}{K\cdot\micro m}.
\end{equation*}
\end{theorem}
\begin{remarque}
  La couleur d'un corps noir n'est pas celle d'un onde monochromatique de longueur d'onde $\lambda_{\rm{max}}$, car le spectre d'un corps noir est relativement large. On pourra comparer ci-dessous la couleur d'un corps noir et de la lumière monochromatique de longueur d'onde $\lambda_{\rm{max}}$.
  \begin{fig}[label]
    \centering
    \includegraphics[width=0.9\columnwidth]{blackbody_color.pdf}
    %\capt{}
  \end{fig}
\end{remarque}

%\begin{remarque}
%La loi de Wien se démontre en dérivant la fonction $u^0_\lambda(\lambda,T)$ donnée par la loi de Planck par rapport à $\lambda$ à $T$ fixée.
%\end{remarque}
% \noindent
% En Fig.~\ref{fig:corps_noir}, chaque courbe est représentée de la couleur du maximum d'émission donné par la loi de Wien.
\begin{exemple}{Couleur des étoiles}
  Utiliser la loi de déplacement de Wien pour répondre aux questions suivantes.
  \begin{quest}
  \item Pourquoi le Soleil apparaît-il jaune  ($T \simeq \SI{5800}{K}$)?
  \item Quelle est la couleur d'une étoile froide ($T \simeq \SI{3500}{K}$) ?
  \item Quelle est la couleur d'une étoile très chaude ($T \simeq \SI{25000}{K}$) ?
  \item Pourquoi n'existe-t-il pas d'étoile verte ?
\end{quest}
\end{exemple}
\examplesol{
% Attention toutefois, la couleur est donnée par l'ensemble du spectre, et pas seulement par le pic d'émission. 
\begin{quest}
  \item	La lumière qui nous provient du Soleil nous apparaît jaune, car $\lambda_{\rm max} \simeq \SI{500}{nm}$ est dans le vert, mais le spectre est suffisamment large pour que l'œil perçoive une couleur jaune (mélange de vert et rouge).
  À température ambiante ($T\simeq\SI{300}{K}$), la majeure partie du spectre se situe dans l'infrarouge, et donc la couleur observée est noire.  
  \item une étoile froide ($T \simeq \SI{3500}{K}$) émet une couleur rouge, car $\lambda_{\rm max} \simeq \SI{830}{nm}$ est dans l'infrarouge proche, mais le spectre est suffisamment large pour que l'œil perçoive une couleur rouge.  
  \item Une étoile très chaude (de température de surface $T \simeq \SI{25000}{K}$) émet une couleur bleue, car $\lambda_{\rm max} \simeq \SI{116}{nm}$ est dans l'ultraviolet, mais le spectre est suffisamment large pour que l'œil perçoive une couleur bleue.
  \item À cause de la sensibilité de l'œil humain et de la largeur spectrale du rayonnement d'équilibre thermique, il n'existe pas d'étoile verte. Le spectre d'une étoile de température intermédiaire (par exemple $T \simeq \SI{6000}{K}$) est suffisamment large pour que l'œil perçoive une couleur blanche, et non verte.
\end{quest}
  }

\parag{Loi de Stefan}
La puissance surfacique $\varphi_0(T) = \int_{0}^{\infty}\dd{\vphi_0}{\lambda}(\lambda,T)\,\d\lambda$ s'obtient par intégration de l'\Eq{planck} sur toutes les longueurs d'onde.
% , soit
% \begin{equation}
%   \varphi^0(T)=\frac{c}{4}u^0(T)=\frac{c}{4}\int_{0}^{\infty}u_{\lambda}^0(\lambda,T)\,\d\lambda.
%   \end{equation}
En calculant l'intégrale\footnote{Poser le changement de variable $u=\frac{hc}{k_{\rm{B}}T\lambda}$ et utiliser $\int_{0}^{\infty}\frac{u^3}{\e^{u}-1}\d u=\frac{\pi^4}{15}$.}, on trouve la loi de Stefan :
\begin{theorem}[Loi de Stefan]\label{form:stefan}
\begin{equation*}
\varphi_0(T)=\sigma T^{4}.
\end{equation*}
\end{theorem}
\noindent La loi de Stefan fait intervenir la constante de Stefan
\begin{equation}
\sigma=\frac{2\pi^5 k_{\rm{B}}^4}{15h^3c^2}=\SI{5,67e-8}{W.m^{-2}.K^{-4}}.
\end{equation}
La puissance surfacique d'équilibre thermique $\varphi_0(T)$, qui représente l'aire sous les courbes de la \Fig{corps_noir}, croît très fortement avec la température (dépendance en $T^4$).


\begin{exemple}{Luminosité des étoiles}
  La luminosité absolue $L$ d'une étoile est la puissance totale qu'elle libère sous forme de rayonnement. Une étoile est de température de surface $T_S$ et est de forme quasi-sphérique de rayon moyen $R$.
  % On donne la constante de Stefan $\sigma = \SI{5.67e-8}{W.m^{-2}.K^{-4}}$.

  %\begin{table}[h!]
  %\renewcommand{\arraystretch}{1.2}
  \begin{center}
  \begin{tabular}{|c|c|c|c|c|}
  \hline
  Étoile & Soleil & Betelgeuse & Deneb & Vega \\
  \hline
   $L$ (W) & $3,8 \cdot 10^{26}$ & $2,4 \cdot 10^{31}$ & $2,3 \cdot 10^{31}$ & $1,4 \cdot 10^{27}$\\
  \hline
  $R$ (km) & $7,0 \cdot 10^5$ & $7,0 \cdot 10^8$ & $1,4 \cdot 10^8$ & $1,8 \cdot 10^6$\\
  \hline
  $T_S$ (K) & $5,8 \cdot 10^3$ & $3,6 \cdot 10^3$ & $8,7 \cdot 10^3$ & $9,6 \cdot 10^3$\\
  % \hline
  % $M$ (kg) & $2,0 \cdot 10^{30}$ & $1,3 \cdot 10^{31}$ & $3,8 \cdot 10^{31}$ & $4,2 \cdot 10^{30}$\\
  \hline
  \end{tabular}
  %\caption{Luminosité, rayon moyen, température de surface et masse de quelques étoiles.}
  %\label{tab:astro}
  \end{center}
  %\end{table}

  	 Expliquer les mesures astronomiques présentées sur le tableau ci-dessus.
\end{exemple}
\examplesol{
La puissance est donnée par la surface de l'étoile multipliée par le puissance  surfacique $L = 4\pi R^2 \vphi$. Si l'étoile se comporte comme une corps noir, la puissance surfacique vérifie la loi de Stefan $\vphi = \vphi_0(T_S) = \sigma T_S^4$.
On vérifie que la luminosité attendue selon la loi de Stefan $L_{\rm att}=4\pi R^2\sigma T_S^4$ s’accorde bien avec la luminosité mesurée $L$ en ordre de grandeur.
\begin{center}
\begin{tabular}{|c|c|c|c|c|}
\hline
Étoile & Soleil & Betelgeuse & Deneb & Vega \\
\hline
 $L$ (W) & $3,8 \cdot 10^{26}$ & $2,4 \cdot 10^{31}$ & $2,3 \cdot 10^{31}$ & $1,4 \cdot 10^{27}$\\
\hline
$L_{\rm att}$ (W) & $4,0 \cdot 10^{26}$ & $5,9 \cdot 10^{31}$ & $8,0 \cdot 10^{31}$ & $2,0 \cdot 10^{28}$\\
\hline
\end{tabular}
\end{center}
La surface des étoiles se comporte donc comme un corps noir en bonne approximation.

%\caption{Luminosité, rayon moyen, température de surface et masse de quelques étoiles.}
%\label{tab:astro}
}


\section{Effet de serre}


\subsection{Bilan radiatif}

\parag{Bilan radiatif}
Soit un corps soumis \`a un rayonnement \'electromagn\'etique quelconque (pas forcément d'équilibre thermique, \Fig{bilan_rad}). Int\'eressons-nous aux différents
\'echanges \'energ\'etiques entre le corps et le rayonnement.  Le rayonnement incident sur le corps, de puissance  $P_i$, peut être redistribu\'e de plusieurs façons.
\begin{liste}
\item Une partie est r\'efl\'echie ou diffus\'ee par la surface de l'objet (puissance $P_r$).
\item Une partie est transmise \`a travers l'objet (puissance
$P_t$).
\item Une partie est absorb\'ee (puissance $P_a$).
\end{liste}
Par ailleurs, le corps peut émettre un rayonnement de puissance surfacique $P_e$. Ces flux surfaciques ne correspondent pas nécessairement à un rayonnement d'équilibre thermique.
La conservation de l'énergie conduit au bilan \'energ\'etique suivant.

\begin{figure}[H]
\begin{center}
\includegraphics[scale=0.7]{schema_flux.pdf}
\end{center}
\caption{Bilan énergétique. Un rayonnement incident de puissance  $P_i$ arrive sur un corps. Une partie est absorbée ($P_a$), une partie transmise ($P_t$), et une partie réfléchie ($P_r$). Le corps peut émettre lui-même un rayonnement de puissance $P_e$ indépendamment du rayonnement incident.}
\label{fig:bilan_rad}
\end{figure}

\begin{theorem}[Bilan enérgétique]\label{form:}
$$
P_i = P_r + P_t + P_a
$$
\end{theorem}
On peut aussi d\'efinir la puissance sortante $P_s = P_r + P_t + P_e$. En régime stationnaire, il y a autant de rayonnement incident sur le corps que sortant, soit $P_i=P_s$ : il s'agit de la condition d'\imp{équilibre radiatif}.
\begin{prop}[Équilibre radiatif]\label{prop:}
En régime stationnaire, la condition d'équilibre radiatif impose $
P_a = P_e$.
\end{prop}
Ainsi, lorsque le corps est à l'équilibre radiatif, tout ce que le corps absorbe est réémis. Attention : l'équilibre radiatif ne préjuge en rien de la nature du rayonnement incident ou sortant, absorbé ou émis (ceux-ci peuvent être à l'équilibre thermique ou non).

\parag{Nature du corps}
On distingue différents types de corps selon leurs réponses à un rayonnement incident, et donc selon la valeur des flux surfaciques transmis, absorbés et réfléchis.
\begin{liste}
\item Un corps est \imp{opaque} si la puissance transmise est nulle,
soit
% \begin{equation}
\eqbox{
  P_t = 0.
  }
% \end{equation}
\item Un corps est \imp{transparent} si la puissance
incidente est int\'egralement transmise, soit
% \begin{equation}
\eqbox{
  P_r = P_a = 0.
  }
% \end{equation}
\item  Un corps partiellement réfléchissant est caractérisé par son \imp{albédo} $A$,  compris entre $0$ et $1$, qui représente la fraction du rayonnement incident qui est réfléchie, soit
% \begin{equation}
\eqbox{
  P_r = A P_i.
  }
% \end{equation}
 % Par exemple un albédo de $\num{0,3}$ signifie que le corps réfléchit $\SI{30}{\percent}$ de l'énergie incidente.
 \item Un corps noir émet un rayonnement d'équilibre thermique. Si le corps noir est de surface émettrice~$S$, on a  d'après la loi de Stefan,
%  \begin{equation}
\eqbox{
  P_e = S \sigma T^4.
  }
%  \end{equation}
\end{liste}



Nous étudions dans cette section des modèles de plus en plus raffinés pour expliquer la température moyenne qui règne à la surface de la Terre. Le rôle de l'atmosphère, au travers de l'effet de serre, en particulier de l'importance de sa concentration en particules absorbantes, sera clarifié.
On regroupe ci-dessous quelques caractéristiques de la Terre et du Soleil, qui est à l'origine du rayonnement incident reçu par la Terre. Dans ce genre d'exercices, il vaut mieux effectuer les \imp{bilans sur les puissances} et non les flux surfaciques, parce que le flux n'est pas distribué de façon uniforme sur la Terre (les rayons sont moins inclinés à l'équateur qu'aux pôles).
%\begin{figure}[h!]
\begin{center}
\begin{tabular}{l|c|c}
 & Terre & Soleil \\
\hline\hline
Rayon & $R_{\rm{T}}=\SI{6,4e6}{m}$ & $R_{\rm{S}}=\SI{7,0e8}{m}$ \\
Température en surface & $T_{\rm{T}} \simeq \SI{15}{\dC}$ & $T_{\rm{S}}=\SI{5,8e3}{K}$ \\
Albédo moyen & $A=\num{0.31}$ & - \\
Distance au Soleil & $d_{\rm{TS}}=\SI{1,5e11}{m}$ & - \\
\end{tabular}
%\caption{Données utiles pour la Terre et le Soleil}
\end{center}
%\end{figure}

\subsection{Modèle sans atmosphère}
\label{sec:serre_sans_atmo}
On néglige pour le moment le rôle de l'atmosphère terrestre. La Terre reçoit un flux incident $\varphi_{S\rightarrow T}$ de la part du Soleil. Une partie de ce flux solaire incident est réfléchi, d'albédo moyen $A =\num{0,31}$, l'autre étant absorbée. Par ailleurs, la Terre est un corps de température $T_T$ et par conséquent elle rayonne également une certaine énergie vers l'espace. On cherche à déterminer dans ce modèle simple la température d'équilibre thermique $T_T$ de la Terre.


\begin{methode}[Étudier l'effet de Serre]
  \item Déterminer la puissance totale $P_S$ émise par le Soleil, assimilé à un corps noir.
  \item En déduire la puissance $P_{\rm S\rightarrow T}$ incidente sur la Terre par des considérations géométriques.
  \item Effectuer un bilan radiatif sur la Terre :
  \begin{liste}
    \item schéma-bilan énergétique pour la Terre, supposée opaque et d'albédo $A$ ;
    %  $P_{\rm S \rightarrow T} = P_r + P_a$ ;
    \item condition d'équilibre radiatif pour la Terre, assimilée à un corps noir.
    %  $P_a=P_e$.
  \end{liste}
  \item Suivant le modèle, effectuer aussi un bilan radiatif sur l'atmosphère.
  \item Conclure en déterminant la température de la surface terrestre $T_T$.
\end{methode}

\begin{figure}[H]
\centering
\subfigure[]{\label{fig:sans_atmo}
\includegraphics[scale=0.6]{schema_terre.pdf}
}\hfill
\subfigure[]{\label{fig:avec_atmo}
\includegraphics[scale=0.6]{schema_terre_atmosphère.pdf}

}
\caption{Bilan d'énergie radiative pour un modèle de Terre : (a) sans atmosphère; (b) avec atmosphère transparente au rayonnement solaire et opaque au rayonnement terrestre.}
\label{fig:bilan_atmo}
\end{figure}
% \underline{Méthode}
% \begin{enumerate}
% \item Exprimer le flux reçu par la Terre $\varphi_{S\rightarrow T}$ en fonction du flux émis par le Soleil $\varphi_S$.
% \item Écrire le bilan radiatif de la Terre en terme de puissances : $P_i = P_r + P_a$.
% \item Hypothèses : équilibre radiatif ($P_e=P_a$) et corps noirs ($\varphi_S = \varphi_0(T_S)$ et $\varphi_e = \varphi_0(T_T)$). En déduire la température de la Terre $T_T$ avec la loi de Stefan.
% \end{enumerate}

\begin{enumerate}
  \item La puissance totale $P_S$ rayonnée par le Soleil a pour expression, d'après la loi de Stefan
  \begin{equation}
  P_{\rm{S}}= \text{(surface du Soleil)} \times \text{(puissance surfacique)} = 4\pi R_{\rm{S}}^2\times \varphi_S = 4\pi R_{\rm{S}}^2 \sigma T_S^4.
  \end{equation}
  \item
  Pour exprimer le flux reçu par la Terre $\varphi_{S\rightarrow T}$ en fonction du flux émis par le Soleil $\varphi_S$, on remarque que la puissance totale $P_S$ rayonnée par le Soleil au niveau de la Terre se répartit sur une sphère de rayon $d_{\rm TS}$. La Terre reçoit donc une puissance surfacique, appelée \imp{constante solaire},
  % \begin{equation}\label{eq:}
  \eqbox[lien_phi]{
    \varphi_{\rm{S}\to\rm{T}} =  \frac{P_{\rm{S}}}{4\pi d_{\rm{TS}}^2} = \sigma T_S^4 \left(\frac{R_{\rm{S}}}{d_{\rm{TS}}}\right)^2 = \SI{1.3e3}{W.m^{-2}}.
    }
  % \end{equation}
  Ce flux n'est pas réparti uniformément sur la surface terrestre à cause de l'inclinaison des rayons par rapport à la surface. La puissance incidente sur la Terre s'obtient en multipliant la puissance surfacique $\varphi_{\rm{S}\to\rm{T}}$ par l'aire $\pi R_T^2$ de la surface projetée perpendiculairement aux rayons solaires, soit
  \begin{equation}\label{eq:P_terre}
  P_{\rm{S}\to\rm{T}}=\pi  R_{\rm{T}}^2 \,\varphi_{\rm{S}\to\rm{T}}.
  %=  \varphi_S \,\pi R_{\rm{T}}^2  \left(\frac{R_{\rm{S}}}{d_{\rm{TS}}}\right)^2.
  \end{equation}
\item
La Terre est opaque : la puissance transmise $P_t = 0$ est nulle. Elle est aussi partiellement réfléchissante d'albédo $A$. Le bilan énergétique appliqué au système Terre se simplifie alors en
\begin{equation}\label{eq:bilan}
P_{S\rightarrow T} = P_{\rm r} + P_{\rm a} = A P_{S\rightarrow T} + P_{\rm a} \Rar P_{\rm a} = (1-A) P_{S \rightarrow T}.
\end{equation}
Cela revient à dire qu'une fraction $1-A$ du rayonnement solaire est absorbé par la Terre, comme représenté sur la \Fig{sans_atmo}.
En supposant que la Terre est à l'équilibre radiatif avec le rayonnement solaire, et en assimilant la Terre à un corps noir, on a
\begin{equation}\label{eq:eq}
P_{\rm{a}} = P_{\rm{e}} = 4\pi R_T^2 \,\sigma T_T^4.
\end{equation}
\item En combinant toutes ces relations, on obtient l'expression de la température de la Terre.
\begin{align*}
(1-A) P_{S \rightarrow T} = (1-A)\pi  R_{\rm{T}}^2 \,\sigma T_S^4 \left(\frac{R_{\rm{S}}}{d_{\rm{TS}}}\right)^2  = 4\pi R_T^2 \,\sigma T_T^4.
\end{align*}
En conclusion,
\eqbox[]{
T_{\rm{T}}=(1-A)^{1/4}\sqrt{\frac{R_{\rm{S}}}{2d_{\rm{TS}}}}\,T_{\rm{S}}.
}
L'application numérique pour la Terre donne $\SI{256}{K}=\SI{-17}{\celsius}$. Cette valeur ne correspond pas à la réalité, puisque le température moyenne mesurée sur Terre vaut 
\eqbox{
T_{\rm{T,mes}}=\SI{288}{K}=\SI{15}{\celsius}.
}
Il faut donc prendre en compte l'atmosphère.
\end{enumerate}


\begin{remarque}
Ce modèle fonctionne bien par contre avec Mercure, qui ne possède pas d'atmosphère. Sa distance au Soleil vaut en moyenne $\SI{5,8e10}{m}$, et son albédo vaut environ $\num{0,068}$. On trouve alors une température de surface moyenne de $\SI{442}{K}=\SI{169}{\celsius}$, en excellent accord avec la réalité ($\SI{440}{K}$ d'après Wikipédia).
\end{remarque}

\subsection{Modèle simple avec atmosphère}
\docu{Composition atmosphérique et bilan radiatif}{Reflets de la Physique}{mars 2013}{}

%Nous allons prendre en compte l'atmosphère.
% L'épaisseur typique de l'atmosphère terrestre est de l'ordre de $e\simeq\SI{30}{km}$. Comme $e\ll R_{\rm{T}}$, on prendra la même surface pour la Terre et pour l'atmosphère, qui l'on assimile à une coquille sphérique de rayon $R_{\rm{T}}$.
Nous prenons l'atmosphère en compte.
Supposons dans ce modèle simple que l'atmosphère est :
\begin{liste}
\item \imp{transparente} au rayonnement solaire : elle transmet intégralement le rayonnement solaire à la Terre;
\item \imp{opaque} au rayonnement terrestre : elle absorbe en totalité le rayonnement issu de la Terre.
\end{liste}

\begin{exemple}{Spectre d'absorption de l'atmosphère}
Commenter ces hypothèses à la lumière des données expérimentales ci-dessous.
  \begin{fig}
  \centering
  \includegraphics[width=0.9\textwidth]{fig1_ED.png}
  % \caption{Spectre d'absorption de l'atmosphère.}
  % \label{fig:asborb}
  \end{fig}
\end{exemple}
\examplesol{
Ces deux hypothèses se justifient parce que les longueurs d'onde des rayonnements émis par le Soleil et la Terre ne sont pas les mêmes. Le Soleil émet d'après la loi de Wien un rayonnement autour de  la longueur d'onde $\lambda_{\rm{S}}\simeq \SI{0,5}{\mu\meter}$ (jaune), alors que la Terre, dont la température de surface est voisine de $\SI{300}{K}$, émet  un rayonnement autour de la longueur d'onde $\lambda_{\rm{T}}\simeq\SI{10}{\mu\meter}$ (infrarouge). L'atmosphère terrestre absorbe peu dans le domaine visible, mais beaucoup dans le domaine IR à cause de la vapeur d'eau. En réalité, on voit que l'atmosphère laisse passer une petite fenêtre de longueur d'onde autour de $\SI{10}{\mu\meter}$. Un modèle plus réaliste prendrait donc en compte une transmission partielle du rayonnement terrestre à travers l'atmosphère.
}

%Pour plus de détails, on se reportera à l'étude documentaire distribuée avec le Chapitre 10


%À partir de ces hypothèses, comment est modifié le modèle précédent ?

Le nouveau bilan radiatif de système composé de l'atmosphère et de la Terre est représenté sur la \Fig{avec_atmo}.
En modélisant la Terre et l'atmosphère par des corps noirs de température respectives $T_{\rm{T}}$ et $T_{\rm{a}}$, la condition d'équilibre radiatif s'écrit pour la Terre d'une part et l'atmosphère d'autre part
\begin{equation}
\begin{cases}
(1-A)P_{\rm{S}\to\rm{T}}+P_{\rm{A}}=P_{\rm{T}},\\
P_{\rm{T}}=2P_{\rm{A}}.
\end{cases}
\end{equation}
En remplaçant par les expressions établies en \Sec{serre_sans_atmo} et en simplifiant, on obtient
\eqbox[]{
T_{\rm{T}}=\JVert{2^{1/4}}(1-A)^{1/4}\sqrt{\frac{R_{\rm{S}}}{2d_{\rm{TS}}}}\,T_{\rm{S}}.
}
On voit que la température est multipliée par un facteur $2^{1/4}$ par rapport au modèle sans atmosphère.
L'application numérique donne $T_{\rm{T}}=\SI{305}{K}=\SI{32}{\celsius}$. Cet effet d'augementation de température est appelé \imp{effet de serre}. Ce modèle conduit toutefois à une température moyenne trop élevée: il est possible de l'améliorer encore en prenant en compte la  transmission partielle du rayonnement terrestre et l'absorption partielle du rayonnement solaire par l'atmosphère.

\begin{fig}[label]
  \centering
  \includegraphics[width=0.7\columnwidth]{Energy_bilan.jpg}
  \capt{Bilan radiatif détaillé de la Terre et de son atmosphère. [Source : \href{https://climate.nasa.gov/faq/19/}{NASA}]}
\end{fig}

\end{document}
